% =========================================================
% Direct Proof
% =========================================================
\paragraph{Direct Proof.}

A direct proof of $P \Rightarrow Q$ assumes $P$ and derives $Q$
through a chain of valid inferences.

\begin{tcolorbox}[colback=propbox, colframe=propborder, arc=2pt,
  left=6pt, right=6pt, top=4pt, bottom=4pt,
  title={\small\textbf{Template: Direct Proof}}, fonttitle=\small\bfseries]
\textbf{Given.} $P$.\\
\textbf{Goal.} $Q$.\\[4pt]
Assume $P$.\\
$\vdots$ \quad [expand definitions; apply lemmas; derive consequences]\\
Therefore $Q$. \qed
\end{tcolorbox}

\begin{remark}[Direct proof is the default --- always try it first]
Direct proof is not one option among equals. It is the default.
Every time you sit down to prove $P \Rightarrow Q$, the first thing you
write is ``Assume $P$'' and the first thing you do is unpack what $P$
says. You use what $P$ gives you, apply definitions and previously proved
results, and derive $Q$.

Reach for contrapositive or contradiction only when direct proof has
genuinely stalled --- when you have unpacked the hypothesis and the goal
and there is no visible path between them. Do not reach for contradiction
just because the statement feels hard. Most statements that feel hard
are hard only because a definition has not been unpacked.
\end{remark}

\begin{remark}[The anatomy of a direct proof]
Every direct proof has three phases that correspond to the first four
steps of the construction algorithm:

\medskip
\noindent \textbf{Phase 1 --- Unpack the hypothesis.} $P$ is a named
concept or a logical statement. Expand it into its operative form.
If $P$ says ``$G$ is a group,'' this licenses associativity, the
existence of an identity, and the existence of inverses. Write those
down as available tools.

\medskip
\noindent \textbf{Phase 2 --- Expand the goal.} $Q$ is also a named
concept. Expand it into what you need to establish. If $Q$ says
``$a = a^{-1}$,'' the operative target is showing $a$ satisfies the
defining property of an inverse: $aa = e$.

\medskip
\noindent \textbf{Phase 3 --- Bridge the gap.} Apply definitions,
lemmas, and axioms to move from the expanded hypothesis to the
expanded conclusion.

\medskip
\noindent When Phase~3 stalls, the cause is almost always that Phase~1
or Phase~2 was not completed. The most common error is having a vague
goal (``show $a = a^{-1}$'') without having written out what that means
in terms of the group axioms.
\end{remark}

\begin{example}[Direct proof in a group]
\textbf{Claim.} In a group $G$, if $a^2 = e$ then $a = a^{-1}$.

\medskip
\noindent\textit{Proof.}
Assume $a^2 = e$, meaning $aa = e$.

We want to show $a = a^{-1}$, which by definition of inverse means
we need $aa = e$ and $ea = a$ (or equivalently, $a$ is its own inverse).

Multiply both sides of $aa = e$ on the right by $a^{-1}$:
\[
(aa)a^{-1} = e \cdot a^{-1}.
\]
By associativity: $a(aa^{-1}) = a^{-1}$.
By the inverse axiom: $ae = a^{-1}$.
By the identity axiom: $a = a^{-1}$.

Hence $a^2 = e$ implies $a = a^{-1}$. \qed

\medskip
\noindent
Each step is justified by exactly one of: the hypothesis ($aa = e$),
the inverse axiom ($aa^{-1} = e$), the identity axiom ($ae = a$),
or associativity. No step is unjustified.
\end{example}
